%=====================================================================
\section{Hydraulics, well-field sizing and performance indices}
\label{sec:hydraulics}

%---------------------------------------------------------------------
\subsection{Pressure drop and pumping power}
\label{sec:hydraulics:dp}

Properties are evaluated at the mean of the inlet and outlet temperatures. With
the flow area $A_{\mathrm{flow}}=\pi r_{i}^{2}$,
\begin{equation}
  u \;=\; \frac{\md}{\rho A_{\mathrm{flow}}},
  \qquad
  \mathrm{Re} \;=\; \frac{\rho\,u\,D_{i}}{\mu},
\end{equation}
and the Darcy friction factor is taken as laminar or Blasius,
\begin{equation}
  f =
  \begin{dcases}
    \dfrac{64}{\mathrm{Re}}, & \mathrm{Re} < 2000, \\[6pt]
    \dfrac{0.3164}{\mathrm{Re}^{1/4}}, & \mathrm{Re} \leq 10^{5},
  \end{dcases}
  \label{eq:friction}
\end{equation}
giving the Darcy--Weisbach frictional drop
\begin{equation}
  \Delta p_{\mathrm{fric}}
  \;=\; f\,\frac{L_{\mathrm{tube}}}{D_{i}}\,\frac{\rho u^{2}}{2} .
  \label{eq:dp-fric}
\end{equation}

Minor losses are added with a lumped coefficient for one hairpin --- a
sharp-edged entrance, two $180^{\circ}$ return bends and an exit:
\begin{equation}
  K_{\mathrm{tot}}
  \;=\; K_{\mathrm{ent}} + 2K_{\mathrm{bend}} + K_{\mathrm{exit}}
  \;=\; 0.5 + 2(2.2) + 1.0 \;=\; 5.9,
  \qquad
  \Delta p_{\mathrm{min}} \;=\; K_{\mathrm{tot}}\,\frac{\rho u^{2}}{2} .
  \label{eq:dp-minor}
\end{equation}
The total drop and the hydraulic power for one tube are then
\begin{equation}
  \Delta p \;=\; \Delta p_{\mathrm{fric}} + \Delta p_{\mathrm{min}},
  \qquad
  \dot W_{\mathrm{pump,tube}} \;=\; \frac{\md\,\Delta p}{\rho} .
  \label{eq:dp}
\end{equation}

The three coefficients in Eq.~\eqref{eq:dp-minor} are hard-coded literals rather
than configuration parameters, and the source describes them as example values.
They are not derived from the completion geometry and cannot be varied from
\code{Case}. At the design point $\Delta p_{\mathrm{min}}$ is small beside
$\Delta p_{\mathrm{fric}}$ because $L_{\mathrm{tube}}/D_{i}\approx 10^{5}$, but
the asymmetry is worth knowing: any sensitivity study on pumping power is
varying only the friction term.
The field total scales by the tube and well counts,
\begin{equation}
  \dot W_{\mathrm{pump}}
  \;=\;
  \dot W_{\mathrm{pump,tube}}\; n_{t}\,\Nwells .
\end{equation}

Two properties of Eq.~\eqref{eq:friction} deserve note. The smooth-tube
assumption sets the roughness to zero, which is optimistic for a well
completion. And no correlation is supplied above $\mathrm{Re}=10^{5}$; the
Blasius branch is simply extrapolated, which is where discharge cases with
elevated flow will land.

\emph{Source:} \code{\_legacy\_physics.calculate\_pressure\_drop}.

%---------------------------------------------------------------------
\subsection{Well-field sizing}
\label{sec:hydraulics:sizing}

\subsubsection*{The v0.1 criterion}

v0.1 sized the field on a single requirement: that the wells absorb the charging
energy within the charging window. Writing $\Delta E_{\mathrm{del}}(\Nwells)$ for
the energy the field takes up over $t_{\mathrm{ch}}$,
\begin{equation}
  \Phi_{\mathrm{heat}}(\Nwells)
  \;\equiv\;
  \frac{\Delta E_{\mathrm{out,HP}}}{\Delta E_{\mathrm{del}}(\Nwells)} - 1
  \;=\; 0 .
  \label{eq:N-heat}
\end{equation}
This is a rate question. Nothing in v0.1 required the wells to \emph{contain}
enough PCM to hold that energy.

\subsubsection*{The capacity criterion}
\label{sec:hydraulics:capacity}

The missing requirement is that the field be able to \emph{contain} the energy
at all, not merely absorb it fast enough:
\begin{equation}
  \boxed{\;
  \Nwells \;=\; \max\bigl(N_{\mathrm{charge\ rate}},\; N_{\mathrm{capacity}}\bigr),
  \qquad
  N_{\mathrm{capacity}}
  = \frac{\Delta E_{\mathrm{out,HP}}}{E_{\mathrm{well}}},
  \;}
  \label{eq:sizing}
\end{equation}
with the capacity of one well written, as in the v0.1 notebook, from latent
\emph{and} sensible heat:
\begin{equation}
  E_{\mathrm{well}}
  = \rhom V_{\mathrm{well}}
    \left[\hm + c_{p,l}\left(T_{3c}-\Tm\right)\right].
  \label{eq:E-well}
\end{equation}
The sensible term is the zero-resistance limit: with no thermal resistance the
liquid would reach the charging inlet temperature. Equation~\eqref{eq:E-well} is
closed form, so unlike the constraint it replaces it costs nothing to evaluate.

This is the quantity the v0.1 notebook computed as \code{N\_wells\_ideal} and
then spent only on costing. In v0.3 its place was taken by a melted-volume
condition, $\varepsilon_{\mathrm{PCM}}(N)=1$, which on inspection was not a
capacity statement at all but ``the well count at which an unthrottled
\SI{10}{\hour} charge exactly saturates the store'' --- a strange thing to size
on, and one that allowed the field to over-charge by \SI{7}{\percent} relative
to the requirement.

Which criterion binds depends on the wall conductivity. With the legacy value
the rate criterion dominates, which is why the omission stayed invisible for so
long. With the correct steel conductivity the rate requirement collapses and
capacity takes over.

\subsubsection*{A reporting rule: $N_{\mathrm{capacity}}$ is a bound, not a result}

Equation~\eqref{eq:sizing} is a maximum, and that has a consequence for how its
parts may be read. $N_{\mathrm{capacity}}$ is the seductive one --- it is closed
form, it is smooth, it responds to every geometric parameter, and \emph{it is
wrong to read on its own}. Whenever the rate criterion binds,
$N_{\mathrm{capacity}}$ moves while the answer does not, and it can move the
opposite way.

The failure is not hypothetical; \S\ref{sec:defects:h3} records a fin sweep in
which $N_{\mathrm{capacity}}$ falls monotonically while $\Nwells$ rises by
\SI{42}{\percent}. Two properties of Eq.~\eqref{eq:E-well} explain it:
$E_{\mathrm{well}} \propto V_{\mathrm{well}}$, and $V_{\mathrm{well}}$ counts
whatever the metal does not occupy --- so $N_{\mathrm{capacity}}\cdot
V_{\mathrm{well}}$ is invariant under any change to the fins, and the column is
pure volume bookkeeping.

The model therefore reports through \code{sizing\_report}, which leads with
$\Nwells$, names the criterion that set it, gives the margin by which the other
is slack, and labels $N_{\mathrm{capacity}}$ as a lower bound. At the design
point the margin is \SI{2.6}{\percent}: the design sits \emph{on} the crossover,
so any parameter sweep switches criterion partway through, and the report warns
when this is so.

\subsubsection*{Why the field is not sized on the discharge}

It is natural to ask whether $\Nwells$ should instead be converged on the
discharge --- the duty that must actually be met --- with the flow converged on
the charge. It cannot, and the reason is structural: \emph{sizing must act on
whichever variable is free.}

During charging the flow is pinned by the specified glide,
$\mdw^{\mathrm{ch}} = \dot Q_{\mathrm{out,HP}}/(c_{p,w}\Delta T_{3C,2C})$, so
$\Nwells$ is the only freedom and the charge can size it. During discharging
$\Nwells$ is already fixed, so the flow is the freedom and the discharge sizes
that instead.

Pinning the discharge flow to its own glide as well --- $\Delta T_{2D,3D}$ is
tied to $\Delta T_{3C,2C}$ by Eq.~\eqref{eq:glide-tied} --- over-determines the
problem, and measurably so:

\begin{center}
\begin{tabular}{@{}r r r@{}}
\toprule
$\Nwells$ & $\varepsilon_{\mathrm{PCM}}$ & $\Delta E_{\mathrm{dc}}/\Delta E_{\mathrm{in,ORC}}$ \\
\midrule
 8    & 1.505 & 0.9285 \\
12.74 & 0.999 & 0.9776 \\
18    & 0.719 & 0.9925 \\
24    & 0.541 & 0.9960 \\
40    & 0.325 & 0.9962 \\
\bottomrule
\end{tabular}
\end{center}

\noindent The delivered energy saturates at $0.996$ and never reaches unity for
\emph{any} well count, because the fluid leaves slightly short of the nominal
outlet temperature and so carries about \SI{0.4}{\percent} less energy per unit
mass than $c_{p}\Delta T$ presumes. There is no root. Letting the discharge flow
float by the required \SI{1.3}{\percent} is precisely what the flow loop exists
to do, and the v0.1 arrangement is therefore correct and is retained.

\subsubsection*{Comparison with the v0.1 ideal well count}
\label{sec:hydraulics:ideal}

The v0.1 notebook printed \code{Ideal number of wells: 12.215} from
\begin{equation}
  m_{\mathrm{ideal}}
  = \frac{\Delta E_{\mathrm{out,HP}}}{\hm + c_{p,l}\left(T_{3c}-\Tm\right)},
  \qquad
  m_{\mathrm{well}} = V_{\mathrm{well}}\,\rho_{l},
  \label{eq:N-ideal-v01}
\end{equation}
which is Eq.~\eqref{eq:E-well} with $\rho_{l}$ in place of $\rhom$. It carries
no thermal resistance and no losses, so it is a genuine \textbf{lower bound}:
the best any design could achieve. Designs with high
$\varepsilon_{\mathrm{PCM}}$ are desirable precisely because they approach it
--- the PCM in the ground is being used rather than merely occupying the
borehole.

\medskip
\noindent\textbf{Correction to version 0.3 of this document.} v0.3 stated that
$N_{\mathrm{inventory}}/N_{\mathrm{ideal}} = 1.0428$ agreed with
$1+\mathrm{Ste} = 1.0474$, and attributed the gap to the sensible heat the model
then omitted. That comparison was made on \emph{inconsistent bases} and the
agreement was a coincidence. Equation~\eqref{eq:N-ideal-v01} uses $\rho_{l}$ and
credits sensible heat; the v0.3 model used $\rho_{s}$ and credited none. On a
common basis:

\begin{center}
\begin{tabular}{@{}l r@{}}
\toprule
variant & $N$ \\
\midrule
Eq.~\eqref{eq:N-ideal-v01} as coded: $\hm+c_{p,l}\Delta T$, $\rho_{l}$ & 12.371 \\
latent only, $\rho_{s}$ & 12.121 \\
$\hm+c_{p,l}\Delta T$, $\rho_{s}$ --- \textbf{consistent with v0.4} & \textbf{11.573} \\
\bottomrule
\end{tabular}
\end{center}

\noindent The two offsetting differences were a factor $0.955$ from the sensible
credit and $0.936$ from the density, net $0.893$. With v0.4 the question
disappears: $N_{\mathrm{capacity}}$ \emph{is} Eq.~\eqref{eq:N-ideal-v01}
evaluated on the model's own conventions, so the comparison is an identity
rather than a check.

\emph{Source:} \code{thums\_core/sizing.py}; reported by \code{sizing\_report}.

\subsubsection*{Discharge flow}

The discharge flow is then found so that the field delivers the ORC's heat
demand over $t_{\mathrm{dc}}$. With $\mathcal{R}=\md^{\mathrm{dc}}_{1}/\md^{\mathrm{ch}}_{1}$
the ratio of per-tube flows,
\begin{equation}
  \frac{\Delta E_{\mathrm{in,ORC}}}
       {\Delta E_{\mathrm{dc}}(\mathcal{R})} - 1 \;=\; 0,
  \label{eq:ratio}
\end{equation}
solved by bisection on $\mathcal{R}$. Under Formulations~B and~C the discharge
march starts from the melt distribution charging left behind, so
Eq.~\eqref{eq:ratio} is bounded by the stored inventory; under Formulation~A it
starts from a bare tube in fully solid PCM and is not.

%---------------------------------------------------------------------
\subsection{Performance indices}
\label{sec:hydraulics:kpi}

The round-trip efficiency, with and without parasitics:
\begin{align}
  \eta_{\mathrm{RTE}}
  &=
  \frac{\bigl(\dot W_{\mathrm{el,out}}
        - \dot W_{\mathrm{pump}}^{\mathrm{dc}}\bigr)\,t_{\mathrm{dc}}}
       {\bigl(\dot W_{\mathrm{el,in}}
        + \dot W_{\mathrm{pump}}^{\mathrm{ch}}\bigr)\,t_{\mathrm{ch}}},
  \label{eq:rte}\\[4pt]
  \eta_{\mathrm{RTE}}^{0}
  &=
  \frac{\dot W_{\mathrm{el,out}}\,t_{\mathrm{dc}}}
       {\dot W_{\mathrm{el,in}}\,t_{\mathrm{ch}}} .
  \label{eq:rte0}
\end{align}
By Eq.~\eqref{eq:budget}, $\eta_{\mathrm{RTE}}^{0}$ depends only on the two cycle
efficiencies, the four electrical and mechanical efficiencies, and $\lambda$. It
is therefore invariant to every storage-model change in this document, and a
useful check that a modification has not leaked into the cycles.

The remaining indices are
\begin{align}
  \varepsilon_{\mathrm{PCM}} &= \frac{\Vmelt n_{t}}{V_{\mathrm{well}}}
    &&\text{melted fraction of inventory}, \\
  E_{\mathrm{well}} &= \frac{\Delta E_{\mathrm{in,ORC}}}{\Nwells}
    &&\text{energy delivered per well}, \\
  \rho_{E} &= \frac{E_{\mathrm{well}}}{V_{\mathrm{well}}}
    &&\text{volumetric energy density}, \\
  f_{\mathrm{pump}} &=
    \frac{\dot W_{\mathrm{pump}}^{\mathrm{ch}}
        + \dot W_{\mathrm{pump}}^{\mathrm{dc}}}{\dot W_{\mathrm{el,out}}}
    &&\text{parasitic fraction}, \\
  \eta_{\mathrm{storage}} &=
    \frac{\Delta E_{\mathrm{discharged}}}{\Delta E_{\mathrm{stored}}}
    &&\text{recovered fraction, Formulations B and C}.
\end{align}

\emph{Source:} \code{thums\_core/kpi.py}; assembled in \code{system.run\_cycle}.

\medskip
\noindent Note that $E_{\mathrm{well}}$ is formed from
$\Delta E_{\mathrm{in,ORC}}$, the energy the ORC consumes, and not from the
energy the store holds. It is a measure of duty per well, not of storage
capacity per well, and the two coincide only when
$\varepsilon_{\mathrm{PCM}}\to 1$.
